\begin{mistake}{2026-06-17}{03}

\begin{problemcontent}
设 \(f(x)\) 在 \((-\infty, +\infty)\) 内是连续的奇函数，\(F(x) = \int_0^{|x|} f(t)\text{d}t\)，则下列选项正确的是（  ）。
A. \(F(x)\) 是不可导的奇函数
B. \(F(x)\) 是可导的偶函数
C. \(F(x)\) 是不可导的偶函数
D. \(F(x)\) 是可导的奇函数
\end{problemcontent}

\begin{solutioncontent}
首先判断奇偶性，计算 \(F(-x)\)：
\[ F(-x) = \int_0^{|-x|} f(t)\text{d}t = \int_0^{|x|} f(t)\text{d}t = F(x) \]
所以 \(F(x)\) 是偶函数，排除选项 A 和 D。

接着讨论可导性。当 \(x \neq 0\) 时，利用链式法则求导：
\[ F'(x) = f(|x|) \cdot (|x|)' = f(|x|) \cdot \frac{x}{|x|} \]
由于 \(f(x)\) 是奇函数：
当 \(x > 0\) 时，\(|x|=x\)，有 \(F'(x) = f(x) \cdot 1 = f(x)\)；
当 \(x < 0\) 时，\(|x|=-x\)，有 \(F'(x) = f(-x) \cdot (-1) = -f(x) \cdot (-1) = f(x)\)。
可见，只要 \(x \neq 0\)，始终有 \(F'(x) = f(x)\)。

最后讨论 \(x = 0\) 处。因为 \(f(x)\) 是定义在 \((-\infty, +\infty)\) 上的连续奇函数，必有 \(f(0) = 0\)。
由于 \(F(x)\) 在 \(x=0\) 处连续，且去心邻域内的导数极限为：
\[ \lim_{x \to 0} F'(x) = \lim_{x \to 0} f(x) = f(0) = 0 \]
由导数极限定理，\(F(x)\) 在 \(x=0\) 处可导，且 \(F'(0) = 0\)。
综上，\(F(x)\) 在 \((-\infty, +\infty)\) 内处处可导。正确答案为 B。
\end{solutioncontent}

\mistakeknowledge{
\begin{itemize}
  \item \textbf{核心公式/定理}：奇偶性定义；导数极限定理（若函数在某点连续，且去心邻域内导数极限存在，则该点可导）。
  \item \textbf{易错点}：看到积分上限带 \(|x|\) 误认为必出现不可导的尖点，未考虑连续奇函数必过原点 \(f(0)=0\) 将左右导数抵消的情形。
  \item \textbf{关联知识}：含有绝对值的复合函数求导公式 \( (|x|)' = \frac{x}{|x|} \quad (x \neq 0) \)。
\end{itemize}}

\end{mistake}
